% Ad Atticum 12.45 --- headnote
% ad-atticum-12-45, 17 May 45 BC, Tusculanum

\textit{Cicero to Atticus, written from the
Tusculanum on the sixteenth day before the Kalends
of June 709 AUC --- 17 May 45 BC (the manuscript
dateline: \textit{Scr.\ in Tusculano xvi K.\ Iun.\
a.\ 709 (45)}). The first letter of the sequence
written from the Tusculanum: Cicero has at last
moved out of Astura, going by way of Lanuvium as he
had announced over the preceding days. The change
of place is also a change of register --- shorter,
more outward-looking, more concerned with the news
from Atticus' household and the literary politics
of the Cato-and-anti-Cato controversy. The opening
note on Attica is brief relief; the report of
Atticus' own \textit{[Greek: ak\=edia]} --- a Greek
term that Cicero will not translate, a listless
indifference that is the elder Atticus' counterpart
to Cicero's grief --- shows that the move closer to
Rome was, in part, for Atticus.}

\textit{The middle sentence is the sequence's
quietest admission about what Astura had been worth:
\textit{ceteroqui ἀνεκτότερα erant Asturae},
``otherwise things were more bearable at Astura.''
The Tusculanum is the right move, but it is not the
easier place; what made Astura possible was its
distance. The aside on Caesar as a neighbour ---
``I would rather have him share a temple with
Quirinus than with Salus'' --- is a sharp dark joke,
all the sharper for being delivered in passing. The
deified Romulus was murdered by the senators of
tradition; the implication, that the dictator might
join him in that fashion, is the kind of thing
Cicero now risks committing to a letter. The close
turns again to Hirtius' anti-Cato pamphlet: spread
it, by all means, so that the praise of the talent
is heard while the theme \textit{[Greek: hypothesis]}
is laughed off.}
